% Part: intuitionistic-logic
% Chapter: tableaux
% Section: introduction

\documentclass[../../../include/open-logic-section]{subfiles}

\begin{document}

\olfileid{int}{tab}{int}

\olsection{પરિચય}

!!^{tableau}s એ !!{signed
  formula}sનાં ચોક્કસ પ્રકારનાં (નીચે તરફ શાખા પાડતાં)
વૃક્ષો છે. ચિહ્નિત સૂત્ર એટલે સત્યમૂલ્ય-ચિહ્ન
($\True$ અથવા $\False$) અને !!a{sentence}ની જોડી:
\[
\sFmla{\True}{!A} \text{ અથવા } \sFmla{\False}{!A}.
\]
!!^a{tableau} કેટલીક \emph{ધારણાઓ}થી શરૂ થાય છે.
પછીનું દરેક !!{signed formula} અનુમાનનો કોઈ નિયમ
લાગુ કરીને મળે છે. કેટલાક નિયમો વૃક્ષની શાખાના છેડે
એક અથવા વધુ !!{signed formula}s ઉમેરે છે; બીજા
નિયમો બે નવા છેડા ઉમેરીને બે શાખાઓ બનાવે છે.
કોઈ નિયમથી મળતા !!{signed formula}sમાંનું
!!{formula} જે !!{signed formula} પર નિયમ લાગુ થયો
તેના સૂત્ર કરતાં ઓછું સંકુલ હોય છે. કોઈ શાખામાં
$\sFmla{\True}{!A}$ અને $\sFmla{\False}{!A}$
બંને હોય, તો તેને \emph{બંધ} કહીએ છીએ.
!!a{tableau}ની દરેક શાખા બંધ હોય, તો આખું
!!{tableau} બંધ છે. બંધ !!{tableau} એવું
!!a{derivation} છે જે બતાવે છે કે !!{tableau}ની
શરૂઆતમાં લેવાયેલા
!!{signed formula}sનો ગણ અસંતોષ્ય છે. આથી
$\Proves$ સંબંધ વ્યાખ્યાયિત કરી શકાય:
$\Gamma \Proves !A$ ત્યારે અને માત્ર ત્યારે જ જ્યારે
એવો કોઈ સાન્ત ગણ~$\Gamma_0 = \{!B_1,
\dots, !B_n\} \subseteq \Gamma$ હોય કે નીચેની
ધારણાઓ માટે બંધ !!{tableau} મળે:
\[
\{\sFmla{\False}{!A}, \sFmla{\True}{!B_1}, \dots, \sFmla{\True}{!B_n}\}.
\]

સ્ફુરણાવાદી તર્ક માટે !!{signed formula}ની સંકલ્પના
વિસ્તૃત કરવી અને સંયોજકોના નિયમો બદલવા પડે છે.
ચિહ્ન ($\True$ અથવા $\False$) ઉપરાંત,
સ્ફુરણાવાદી !!{tableau}sનાં !!{formula}sને
\emph{ઉપસર્ગો}~$\sigma$ પણ હોય છે.
\sourcecorrection{OLMD-097}{મૂળના આ વાક્યમાં
``મોડલ ટેબ્લો'' છે, પરંતુ અહીંનું પ્રકરણ અને તેના
નિયમો સ્ફુરણાવાદી ટેબ્લો વિશે છે; વિષયનું નામ
સુસંગત કર્યું છે.}
ઉપસર્ગો ધન પૂર્ણાંકોની ખાલી ન હોય તેવી શ્રેણીઓ છે,
એટલે $\sigma \in (\PosInt)^* \setminus \{\emptyseq\}$.
આવા ઉપસર્ગો લખતી વખતે આપણે આસપાસના
$\tuple{\ }$ છોડીએ છીએ અને જુદા !!{element}s
વચ્ચે $,$ને બદલે~$.$ મૂકીએ છીએ. જો $\sigma$
ઉપસર્ગ હોય, તો $\sigma.n$ એ $\sigma \concat \tuple{n}$
છે; દાખલા તરીકે, $\sigma = 1.2.1$ હોય, તો
$\sigma.3$ એ $1.2.1.3$ છે. આમ, દાખલા તરીકે,
\[
\sFmla{\True}{!A \lif (!B \lif !C)}[1.2]
\]
એ \emph{ઉપસર્ગવાળું !!{signed formula}} છે
(અથવા ટૂંકમાં \emph{ઉપસર્ગવાળું !!{formula}}).

સહજ રીતે, ઉપસર્ગ નિદર્શમાં એવું જગત સૂચવે છે જે
!!a{tableau}ની કોઈ શાખા પરનાં !!{formula}sને
સંતોષી શકે; અને $\sigma$ કોઈ જગતનું નામ હોય, તો
$\sigma.n$ એ~$\sigma$ નામના જગતમાંથી સુલભ જગતનું
નામ આપે છે.

સ્ફુરણાવાદી નિદર્શોમાં સુલભતા સંબંધ સ્વવાચક અને
પરંપરિત હોય છે. ઉપસર્ગોની દૃષ્ટિએ તેનો અર્થ એ કે
$\sigma$ પોતે~$\sigma$માંથી સુલભ છે, અને
$\sigma$ને વિસ્તારીને મળતો કોઈ પણ ઉપસર્ગ,
એટલે $\sigma.n_1.\cdots.n_k$ સ્વરૂપનો ઉપસર્ગ,
પણ સુલભ છે. $\sigma$ પોતે અને તેનો કોઈ પણ
વિસ્તાર દર્શાવવા $\sigma.*$ સંજ્ઞા લઈએ.
બીજા શબ્દોમાં, $\sigma.*$માં બરાબર એ બધા
ઉપસર્ગો છે જે~$\sigma$માંથી સુલભ છે.

\end{document}
